% simulant — Cover letter (ACS Sensors), 2026-11-02 신설.
%
% JHM 판본(paper/cover_letter.tex, 2,417단어)을 줄여 쓴 것이 아니라 새로 쓴다.
% 이유는 두 가지다.
%
% (1) 순서가 다르다. JHM 쪽은 지정문서 기록으로 열었다. ACS Sensors 독자에게
%     헤드라인은 응답공간 결과이므로 그쪽으로 연다.
% (2) 스코프 관문을 정면으로 답해야 한다. ACS Sensors 저자 안내(2026-10-08
%     실측)는 복잡 시료에서의 센서 실증과 기존 분석법과의 상관을 요구하고,
%     '알려진 센싱 개념'에 대한 심사없는 반려 조항을 둔다. 이 원고는 센서
%     개발 논문이 아니므로 셋 다 해당하지 않는다. 숨기면 데스크 리젝이고,
%     먼저 꺼내 왜 그래도 스코프 안인지 말하는 것이 유일한 길이다.
%
% em-dash 금지 / AI 문체 금지 규칙 적용.
\documentclass[11pt,letterpaper]{article}
\usepackage[margin=2.5cm]{geometry}
\usepackage{times}
\usepackage[hidelinks]{hyperref}
\setlength{\parindent}{0pt}
\setlength{\parskip}{8pt}
\pagestyle{empty}

\begin{document}

{\small
Chemical, Biological, Radiological and Nuclear (CBRN) Defense Research
Institute\\
Ministry of National Defense, Seoul, Republic of Korea\\
2 November 2026}

\vspace{6pt}

% [확인 필요] 투고 시스템에서 배정된 Associate Editor 성함을 확인하고
%            "Dear Dr. <Name>," 로 교체할 것. ACS Sensors 는 복수
%            Associate Editor 체제다.
Dear Editors,

We submit \textit{Descriptor Similarity Does Not Predict Array Response:
Chemical Warfare Agent Simulants Tested Against Sixteen Authentic Agents}
for consideration as an Article.

\textbf{What the paper is about, and why we bring it here.} Chemical warfare
agent detection is almost always developed and qualified on simulants, because
the agents themselves are available to few laboratories. Every such
qualification rests on one assumption: that a compound chosen to resemble an
agent on molecular descriptors will resemble it in what the detector measures.
We tested that assumption directly, against colorimetric response measured on
sixteen authentic agents spanning nerve, blood, choking, vesicant and
A-series compounds. This journal sets the standard for whether a detection
claim has been demonstrated against the analyte it names. The question we ask
is one step earlier and bears on every simulant-qualified device in the
literature: whether the surrogate analyte was the right one to begin with.

\textbf{Three counts state the result.} Across the sixteen agents, exactly
one designated simulant has ever been measured beside them, and that single
co-measurement sits in one run of one array. Twelve exact Mantel tests and six
Fisher combinations of them leave no agreement between descriptor distance and
response distance below $q = 0.05$. On the thirty-five-compound panel, no agent
of eleven has a valid match to an incumbent simulant, at a
Benjamini--Hochberg adjusted value of 0.6173. The array itself is not in
question: it classifies these agents at a cross-validated accuracy of 0.964 on
the seven-agent run and 1.000 on the Novichok run. What fails is the transfer
from descriptor space to response space.

\textbf{The specific case will interest sensor developers.} Dimethyl
methylphosphonate is the standard G-agent simulant and the benchmark compound
for decontamination efficiency. Measured beside the agents in the ultraviolet
channel, its response pattern correlates at $+0.76$ with tabun and $+0.65$ with
sarin but at $-0.49$ with soman, $-0.72$ with cyclosarin and $-0.70$ with VX.
Its similarity to sarin also falls with concentration, from $+0.77$ at
500~$\mu$M to $+0.36$ at 10~$\mu$M, which is the end at which a field detector
operates. A surrogate that resembles the agent only at high dose is of limited
use to the measurement it stands in for.

\textbf{What this paper is not, stated plainly.} It reports no new sensor, no
demonstration in a complex matrix, and no correlation against an established
analytical method, and we recognise that your guidelines ask for those of a
sensing paper. This is not a sensing paper. It is a measurement-science audit
of the surrogate on which sensing papers are qualified, and the data it audits
could not be produced any other way: the agents cannot be put into a complex
environmental matrix outside a designated laboratory, and the comparison it
makes is between two descriptions of the same compounds rather than between a
sensor and a reference method. We would rather you weigh that against your
scope now than have it surface in review.

\textbf{Data provenance.} The colorimetric response data are reanalysed from
three earlier studies by this group, in \textit{Journal of Hazardous
Materials}, \textit{Discover Sensors} and \textit{Sensors and Actuators B},
all cited in the manuscript. Those papers reported classification performance
for three separate agent sets. None of them compared response to molecular
descriptors, none placed the three sets in one analysis, and no figure, table
or passage is reproduced here. The descriptor calculations, the permutation
tests and every result in this manuscript are new.

\textbf{What we propose instead of a replacement compound.} We propose none,
because selecting one on descriptor similarity would repeat the error the
paper is about. What the manuscript specifies is the measurement that would
settle the question: a bridging run in which candidate simulants and the agents
they stand in for are exposed on the same plate, in the same optical channels,
at the same concentrations, within one session. That run yields a response-space
distance between a simulant and its agent measured under identical conditions,
which no amount of descriptor work can supply. It needs no new chemistry and
no new instrument.

\textbf{Reporting and availability.} All agents were produced and handled at an
Organisation for the Prohibition of Chemical Weapons designated laboratory
under Chemical Weapons Convention procedures. Analysis code, descriptor tables,
parsed regulatory values with their sources and every figure script are
deposited in a public repository cited in the manuscript, and all permutation
tests are exhaustive rather than sampled, so every number reported can be
reproduced from the deposit. The manuscript has not been published elsewhere
and is not under consideration at another journal. The authors declare no
competing financial interest.

\vspace{6pt}

Yours sincerely,

\vspace{14pt}

Ku Kang, on behalf of the authors\\
{\small Corresponding author: Yoon Jeong Jang, yjjangchem@korea.kr}

\end{document}
